\section{ТЕСТИРОВАНИЕ И АНАЛИЗ РЕЗУЛЬТАТОВ}

\subsection{Методика тестирования и тестовое окружение}

Тестирование разработанной системы проводилось на двух независимых наборах тестовых данных, охватывающих различные предметные области. Первый набор основан на схеме Sakila~\cite{sakila_db} --- эталонной демонстрационной базе данных, разработанной командой MySQL. Второй --- Brazilian E-Commerce Public Dataset by Olist~\cite{olist_dataset}, публичный набор данных электронной торговли с платформы Kaggle, содержащий сведения о реальных заказах.

Каждый набор данных был развёрнут в трёх экземплярах СУБД: MySQL~8, PostgreSQL~15 и MariaDB~10.x --- в Docker-контейнерах. Серверная часть системы запускалась посредством Uvicorn (порт~8000), клиентская --- через Next.js-сервер (порт~3000). База данных истории тестов (\texttt{project-db}) развёрнута в контейнере PostgreSQL~15 на порту~5433. Для схемы Sakila подключения работали на портах~3306 (MySQL), 5437 (PostgreSQL) и~3307 (MariaDB). Для набора Olist использовались порты~3308 (MySQL), 5434 (PostgreSQL) и~3309 (MariaDB).

Для совместного тестирования трёх экземпляров СУБД как единого объекта применялась функция \textbf{логической базы данных} --- именованной группы физических подключений под единым профилем схемы. Система автоматически определила профиль \texttt{sakila\_like} для схемы Sakila и \texttt{olist\_like} для набора Olist, после чего сгенерировала пакеты запросов, специфичные для каждого профиля.

Тестирование выполнялось в двух форматах. В формате \textbf{одиночного прогона} (режим \texttt{per\_test}) все три подключения тестировались одновременно при фиксированных параметрах нагрузки; результаты использовались для попарного сравнения СУБД. В формате \textbf{серийного тестирования} (режим \texttt{series}) выполнялось несколько прогонов с различными уровнями нагрузки для анализа масштабируемости. Параметры нагрузки задавались числом виртуальных пользователей~(VU) и числом итераций на пользователя; время прогрева во всех прогонах составляло~5~секунд.

Статистическая значимость различий оценивалась критерием Манна--Уитни (уровень значимости $\alpha = 0{,}05$) с поправкой Бенджамини--Хохберга (FDR) для множественных сравнений. Практическая значимость определялась по размеру эффекта Коэна~$d$: $d < 0{,}2$ --- пренебрежимый, $0{,}2 \leq d < 0{,}5$ --- малый, $0{,}5 \leq d < 0{,}8$ --- средний, $d \geq 0{,}8$ --- большой.

Сравнительные характеристики тестовых наборов данных приведены в таблице~\ref{tab:test_databases}.

\begin{table}
    \centering
    \caption{Характеристики тестовых наборов данных}
    \label{tab:test_databases}
    \begin{tabularx}{\textwidth}{|>{\RaggedRight}X|>{\RaggedRight}X|>{\RaggedRight}X|>{\RaggedRight}X|}
        \hline
        Набор данных & Профиль схемы & Таблиц & Тип первичного ключа \\
        \hline
        Sakila (MySQL, порт~3306) & \texttt{sakila\_like} & 16 & INT \\
        \hline
        Pagila (PostgreSQL, порт~5437) & \texttt{sakila\_like} & 16 & INT \\
        \hline
        Makila (MariaDB, порт~3307) & \texttt{sakila\_like} & 16 & INT \\
        \hline
        Brazilian E-com (PostgreSQL, порт~5434) & \texttt{olist\_like} & 9 & CHAR(32) \\
        \hline
        Brazilian E-com (MySQL, порт~3308) & \texttt{olist\_like} & 9 & CHAR(32) \\
        \hline
        Brazilian E-com (MariaDB, порт~3309) & \texttt{olist\_like} & 9 & CHAR(32) \\
        \hline
    \end{tabularx}
\end{table}

Схема Sakila содержит 16~таблиц с целочисленными первичными ключами, что обеспечивает высокую скорость поиска по индексам. Набор Olist состоит из~9~таблиц с первичными ключами типа \texttt{CHAR(32)} в формате MD5-хэша и содержит около~100~000~записей заказов, что приводит к более высокой латентности запросов по сравнению с Sakila. Данные различия в характеристиках двух наборов позволяют исследовать поведение СУБД в принципиально разных условиях нагрузки.

\subsection{Тестирование на наборе данных Sakila}

\subsubsection{Сравнение СУБД при смешанной нагрузке}

Первый прогон на наборе Sakila выполнялся в режиме \texttt{per\_test} с параметрами: 10 виртуальных пользователей, 50 итераций на пользователя, сценарий \texttt{mixed\_light}. Автоматически сгенерированный пакет запросов (\texttt{mixed\_light::Sakila::common}) содержал пять операций на таблице \texttt{payment}: выборку строки по первичному ключу (вес~30), обновление временной метки (вес~20), выборку по внешнему ключу \texttt{customer\_id} (вес~15), JOIN-запрос с таблицей \texttt{customer} (вес~20) и инкремент поля \texttt{amount} (вес~20). Всего выполнено~1~500 транзакций за~29,8~с.

Метрики производительности для каждой СУБД приведены в таблице~\ref{tab:sakila_mixed_metrics}.

\begin{table}
    \centering
    \caption{Метрики производительности (Sakila, mixed\_light, 10~VU / 50~итераций)}
    \label{tab:sakila_mixed_metrics}
    \begin{tabularx}{\textwidth}{|>{\RaggedRight}X|>{\RaggedRight}X|>{\RaggedRight}X|>{\RaggedRight}X|>{\RaggedRight}X|}
        \hline
        СУБД & Среднее, мс & p95, мс & p99, мс & TPS, зап/с \\
        \hline
        PostgreSQL (Pagila) & 12{,}11 & 10{,}55 & 280{,}0 & 500 \\
        \hline
        MySQL (Sakila) & 11{,}40 & 25{,}06 & 68{,}45 & 250 \\
        \hline
        MariaDB (Makila) & 8{,}25 & 10{,}71 & 59{,}08 & 500 \\
        \hline
    \end{tabularx}
\end{table}

Все три СУБД завершили тест без ошибок (error\_rate~=~0). MariaDB продемонстрировала наименьшую среднюю задержку~--- 8,25~мс против 11,40~мс у MySQL и 12,11~мс у PostgreSQL. PostgreSQL показал лучший показатель p95 (10,55~мс), тогда как его p99 достиг 280,0~мс~--- аномально высокое значение, указывающее на наличие редких выбросов. MySQL показал наихудший p95 (25,06~мс) при наименьшей пропускной способности (250~зап/с).

Попарный анализ выявил статистически значимое различие между PostgreSQL и MariaDB: MariaDB быстрее на~31,9\% (критерий Манна--Уитни, $p < 0{,}001$, после FDR-коррекции), однако размер эффекта пренебрежимо мал ($d = 0{,}14$). Различие между MySQL и MariaDB статистически незначимо ($p = 0{,}059$, $d = 0{,}34$~--- малый эффект). Ресурсные метрики демонстрируют равномерную загрузку: CPU~--- 1,4--5,1\%, коэффициент попаданий в кэш~--- 99,04--99,86\%, взаимных блокировок не зафиксировано.

Таким образом, при умеренной смешанной нагрузке все три СУБД работают корректно, а различия в производительности статистически незначительны с точки зрения практического эффекта.

\subsubsection{Масштабируемость при росте нагрузки}

Для анализа масштабируемости проведено серийное тестирование в режиме \texttt{series}: два прогона сценария \texttt{read\_only} на том же наборе Sakila с уровнями нагрузки 5~VU~/ 30~итераций (базовый) и 30~VU~/ 30~итераций. Пакет запросов (\texttt{read\_only::Sakila::common}) включал шесть SELECT-операций: выборку строки по PK, выборку нескольких столбцов по PK, выборку по FK, JOIN с одной таблицей, JOIN с двумя таблицами и агрегацию \texttt{COUNT(*)}.

Траектории средней задержки и пропускной способности по уровням нагрузки приведены в таблице~\ref{tab:sakila_series}.

\begin{table}
    \centering
    \caption{Траектории метрик по уровням нагрузки (Sakila, read\_only)}
    \label{tab:sakila_series}
    \begin{tabularx}{\textwidth}{|>{\RaggedRight}X|>{\RaggedRight}X|>{\RaggedRight}X|>{\RaggedRight}X|>{\RaggedRight}X|}
        \hline
        СУБД & \multicolumn{2}{c|}{Среднее, мс} & \multicolumn{2}{c|}{TPS, зап/с} \\
        \cline{2-5}
        & 5~VU / 30~ит. & 30~VU / 30~ит. & 5~VU / 30~ит. & 30~VU / 30~ит. \\
        \hline
        PostgreSQL (Pagila) & 10{,}12 & 52{,}94 & 415{,}7 & 692{,}1 \\
        \hline
        MySQL (Sakila) & 7{,}58 & 26{,}56 & 150{,}0 & 178{,}0 \\
        \hline
        MariaDB (Makila) & 9{,}01 & 25{,}71 & 150{,}0 & 781{,}0 \\
        \hline
    \end{tabularx}
\end{table}

При переходе от 5 к 30~VU поведение СУБД принципиально различается. MariaDB демонстрирует наибольшую эластичность: пропускная способность выросла с~150 до~781~зап/с (+421\%), при этом деградация p95 составила лишь~+59\%, а p99 снизился на~21,6\%~--- то есть хвост распределения стал более предсказуемым. Коэффициент эластичности MariaDB~--- 0,84~--- существенно выше, чем у PostgreSQL (0,13) и MySQL (0,04).

PostgreSQL показал наибольшую деградацию хвоста: p95 вырос на~638\%, p99~--- на~657\% (с~10,1 до~76,7~мс). Рост средней задержки MySQL статистически подтверждён: +250,4\% ($d = 2{,}94$, большой эффект, $p < 0{,}001$, FDR-скорректировано). Ни у одной из СУБД в данной серии не зафиксировано ошибок выполнения запросов.

Блок \texttt{parameter\_impacts} аналитического модуля системы зафиксировал ключевые изменения при увеличении VU с~5 до~30: у Pagila~--- рост средней задержки на~422,9\% и p99 на~657,3\%; у MySQL~--- рост средней задержки на~250,3\% при одновременном снижении коэффициента вариации (CV) на~68,2\%, что указывает на улучшение стабильности распределения при высокой нагрузке. Полученные результаты подтверждают, что система корректно фиксирует разнонаправленные эффекты масштабирования для различных СУБД.

\subsection{Тестирование на наборе Brazilian E-Commerce}

\subsubsection{Сравнение СУБД при смешанной нагрузке}

Первый прогон на наборе Brazilian E-Commerce выполнялся в режиме \texttt{per\_test} с параметрами: 10~VU, 50~итераций, сценарий \texttt{mixed\_light}. Пакет запросов (\texttt{mixed\_light::Brazilian E-com::common}) содержал пять операций: инкремент поля по PK в таблице \texttt{olist\_customers\_dataset} (вес~20), выборку строки по PK из той же таблицы (вес~30), вставку строки в \texttt{olist\_geolocation\_dataset} (вес~15), выборку строк по FK из \texttt{olist\_order\_items\_dataset} (вес~15) и JOIN-запрос \texttt{olist\_orders\_dataset} с \texttt{olist\_customers\_dataset} (вес~20). Всего выполнено~1~500 транзакций за~122,7~с.

Метрики производительности для каждой СУБД приведены в таблице~\ref{tab:ecom_mixed_metrics}.

\begin{table}
    \centering
    \caption{Метрики производительности (Brazilian E-com, mixed\_light, 10~VU / 50~итераций)}
    \label{tab:ecom_mixed_metrics}
    \begin{tabularx}{\textwidth}{|>{\RaggedRight}X|>{\RaggedRight}X|>{\RaggedRight}X|>{\RaggedRight}X|>{\RaggedRight}X|}
        \hline
        СУБД & Среднее, мс & p95, мс & p99, мс & TPS, зап/с \\
        \hline
        PostgreSQL & 12{,}94 & 10{,}46 & 293{,}87 & 250 \\
        \hline
        MySQL & 11{,}34$^{*}$ & 11{,}34$^{*}$ & 11{,}34$^{*}$ & 676 \\
        \hline
        MariaDB & 7{,}75 & 12{,}14 & 63{,}45 & 500 \\
        \hline
    \end{tabularx}
    {\footnotesize $^{*}$~Для MySQL доступна только~1 выборка латентности; статистический тест по данной метрике невозможен.}
\end{table}

По результатам прогона MySQL показал наибольшую пропускную способность. MariaDB продемонстрировала наименьшую среднюю задержку~--- 7,75~мс; статистически значимое различие зафиксировано только в паре PostgreSQL~--- MariaDB: MariaDB быстрее на~40,1\% ($p < 0{,}001$, FDR-скорректировано, $d = 0{,}17$~--- пренебрежимый эффект).

Примечательно, что коэффициент попаданий в кэш буферного пула PostgreSQL составил~92,21\%~--- существенно ниже, чем у MySQL (99,93\%) и MariaDB (99,48\%). Это объясняется тем, что объём набора Olist превышает рабочий объём буферного пула контейнера PostgreSQL, и значительная часть операций чтения требует обращения к диску. Ни одной ошибки выполнения запросов зафиксировано не было.

\subsubsection{Транзакционная нагрузка (OLTP)}

Второй прогон на наборе Brazilian E-Commerce проводился в режиме \texttt{per\_test} с параметрами: 15~VU, 50~итераций, сценарий \texttt{oltp}. Пакет запросов (\texttt{oltp::Brazilian E-com::common}) содержал четыре операции: выборку строки по PK и проекционную выборку нескольких столбцов из \texttt{olist\_customers\_dataset} (веса~30 и~20), выборку позиций заказа из \texttt{olist\_order\_items\_dataset} (вес~15) и обновление поля записи заказа (вес~35). Всего выполнено~2~250 транзакций; длительность прогона составила~24,1~с для PostgreSQL и~15,8--16,3~с для MariaDB и MySQL.

Метрики производительности при транзакционной нагрузке приведены в таблице~\ref{tab:ecom_oltp_metrics}.

\begin{table}
    \centering
    \caption{Метрики производительности (Brazilian E-com, OLTP, 15~VU / 50~итераций)}
    \label{tab:ecom_oltp_metrics}
    \begin{tabularx}{\textwidth}{|>{\RaggedRight}X|>{\RaggedRight}X|>{\RaggedRight}X|>{\RaggedRight}X|>{\RaggedRight}X|}
        \hline
        СУБД & Среднее, мс & p95, мс & p99, мс & TPS, зап/с \\
        \hline
        PostgreSQL & 32{,}15 & 23{,}02 & 786{,}54 & 375 \\
        \hline
        MySQL & 32{,}73 & 112{,}59 & 143{,}80 & 53 \\
        \hline
        MariaDB & 21{,}12 & 25{,}81 & 83{,}26 & 375 \\
        \hline
    \end{tabularx}
\end{table}

При транзакционном сценарии MariaDB демонстрирует наилучшие показатели по средней задержке (21,12~мс) и p95 (25,81~мс). Различие MariaDB и MySQL статистически значимо: MariaDB быстрее на~35,5\% ($p = 0{,}020$, $d = 0{,}78$~--- средний эффект). Различие MariaDB и PostgreSQL также значимо: MariaDB быстрее на~34,3\% ($p < 0{,}001$, $d = 0{,}14$~--- пренебрежимый эффект).

Наиболее примечательная аномалия~--- пропускная способность MySQL (53~зап/с) в~7 раз ниже, чем у PostgreSQL и MariaDB (375~зап/с каждая). Вместе с тем p99 MySQL (143,80~мс) существенно лучше, чем у PostgreSQL (786,54~мс). Таким образом, MySQL при транзакционной нагрузке на данном наборе данных обрабатывает меньший поток запросов, но не демонстрирует экстремальных выбросов хвоста, в отличие от PostgreSQL. Все три СУБД завершили прогон без ошибок.

\subsubsection{Масштабируемость при росте нагрузки}

Для набора Brazilian E-Commerce серийное тестирование проводилось со сценарием \texttt{read\_only} на двух уровнях нагрузки: 5~VU~/ 40~итераций и 30~VU~/ 40~итераций. Пакет запросов (\texttt{read\_only::Brazilian E-com::common}) включал шесть SELECT-операций на таблицах \texttt{olist\_customers\_dataset}, \texttt{olist\_order\_items\_dataset} и \texttt{olist\_orders\_dataset}, включая JOIN-запросы через FK-связи и агрегацию.

Траектории метрик по уровням нагрузки приведены в таблице~\ref{tab:ecom_series}.

\begin{table}
    \centering
    \caption{Траектории метрик по уровням нагрузки (Brazilian E-com, read\_only)}
    \label{tab:ecom_series}
    \begin{tabularx}{\textwidth}{|>{\RaggedRight}X|>{\RaggedRight}X|>{\RaggedRight}X|>{\RaggedRight}X|>{\RaggedRight}X|}
        \hline
        СУБД & \multicolumn{2}{c|}{Среднее, мс} & \multicolumn{2}{c|}{TPS, зап/с} \\
        \cline{2-5}
        & 5~VU / 40~ит. & 30~VU / 40~ит. & 5~VU / 40~ит. & 30~VU / 40~ит. \\
        \hline
        PostgreSQL & 693 & 3~426 & 6{,}0 & 6{,}7 \\
        \hline
        MySQL & 1~292 & 6~428 & 4{,}5 & 4{,}3 \\
        \hline
        MariaDB & 1~366 & 6~769 & 4{,}8 & 4{,}7 \\
        \hline
    \end{tabularx}
\end{table}

Результаты серийного тестирования на наборе Olist демонстрируют принципиально иное поведение по сравнению с Sakila: при шестикратном росте числа VU пропускная способность всех трёх СУБД остаётся практически неизменной (изменение в пределах~2--13\%, статистически незначимо). Модуль анализа зафиксировал достижение точки насыщения (\texttt{saturation\_point}) для MySQL и MariaDB на уровне~30~VU; у PostgreSQL точка насыщения не определена формально, однако прирост TPS незначителен (+13\%). Коэффициенты эластичности отрицательны или близки к нулю: MySQL~--- $-$0,012, MariaDB~--- $-$0,001, PostgreSQL~--- 0,026.

Средняя задержка при этом растёт значительно: у MySQL на~397\% ($d = 1{,}02$, большой эффект), у MariaDB на~396\% ($d = 0{,}95$, большой эффект), у PostgreSQL на~394\% ($d = 1{,}03$, большой эффект). Деградация p99 составила: PostgreSQL~--- +264\%, MySQL~--- +319\%, MariaDB~--- +409\%. При базовой нагрузке PostgreSQL показывает наименьшую среднюю задержку (693~мс против 1~292~мс у MySQL и 1~366~мс у MariaDB), что объясняется более эффективным планировщиком запросов при работе с CHAR(32)-ключами. Во всех прогонах ошибок выполнения не зафиксировано.
